\chapter{El término afín del ruido}
\label{anexo:factor_negativo}

% =====================================================================================
% ANEXO CORTO A PROPOSITO. Sale del capitulo 5 para no desviar la linea de trabajo: alli
% §5.5 solo declara la simplificacion y remite aqui, y §5.3.2 apunta aqui al hablar de los
% factores negativos.
%
% 🔴 ES EL UNICO SITIO DE LA MEMORIA donde se desarrolla el modelo afin. No reponer la
%    formula ni el descarte de `b` en §5.5: quedaria dicho dos veces.
%
% Fuente de las cifras: doc/changes/2026-08_cambio_decisiones.md §1 y §3 (dataset v2,
% ajuste sobre una submuestra de 3.000 circuitos). Todas citables.
%
% ⚠️ El paso de `f = f_real + b/e` es algebra, no una medicion: sale de dividir la
%    ecuacion afin por `e`. No hace falta respaldarlo con ninguna cifra.
% =====================================================================================

El factor de supervivencia que define el capítulo~\ref{cap:variable} asume que el ruido es
\textbf{multiplicativo puro}: que el valor medido es el exacto multiplicado por una fracción
que sobrevive. Medido, no lo es del todo. El modelo real es \textbf{afín},
\begin{equation}
\label{eq:modelo_afin}
\langle O \rangle_{\text{ruidoso}} = f \cdot \langle O \rangle_{\text{exacto}} + b ,
\end{equation}
\noindent donde $b$ es un desplazamiento que \textbf{no depende de cuánta señal había}. Tiene
explicación física: la relajación hace caer los qubits hacia $\ket{0}$, cuyo valor de $Z$ es
$+1$, así que el ruido empuja la base $Z$ hacia valores positivos aunque el valor exacto sea
cero.

\section{Por qué el factor se dispara, y por qué cambia de signo}

La consecuencia se ve dividiendo la ecuación~\eqref{eq:modelo_afin} por el valor exacto, que
es lo que hace la definición del factor. Escribiendo $e$ para el valor exacto:
\begin{equation}
\label{eq:factor_afin}
\frac{\langle O \rangle_{\text{ruidoso}}}{e} = f + \frac{b}{e} .
\end{equation}

\noindent El factor que se mide no es $f$, sino $f$ \textbf{más un término que crece cuando el
denominador se encoge}. Con un valor exacto grande, $b/e$ es despreciable y las dos
formulaciones coinciden. Con un valor exacto diminuto (que es el caso típico en este conjunto
de datos, donde la mediana del valor absoluto ronda las tres centésimas) ese término domina, y
si además $b$ tiene el signo contrario al de $e$, \textbf{el cociente entero se vuelve
negativo}. De ahí salen los factores negativos del apartado~\ref{sec:umbral}.

\section{Por qué no se modela}

Para simplificar el trabajo.